\begin{proof}[Proof of \cref{obj:lem:classical-legendre-facts}]
Write the shifted Legendre polynomials in the normalization
\[
R_j(x):=\sum_{m=0}^{j}(-1)^m
  \binom{j}{m}\binom{j+m}{m}x^m,
\qquad
P_j^{\mathrm{Leg}}(t)=R_j\!\left(\frac{1-t}{2}\right),
\qquad j\in\mathbb N,\quad x,t\in\mathbb R.
\]

\begin{enumerate}
\item \emph{Orthogonality.}
The affine substitution \(t=1-2x\), including reversal of the integration endpoints, gives
\[
\int_0^1R_j(x)R_k(x)\,dx
=-\frac12\int_1^{-1}
  P_j^{\mathrm{Leg}}(t)P_k^{\mathrm{Leg}}(t)\,dt
=\frac12\int_{-1}^1
  P_j^{\mathrm{Leg}}(t)P_k^{\mathrm{Leg}}(t)\,dt.
\]
We therefore compute the inner product on the unit interval.

For \(k\in\mathbb N\), define
\[
\mathcal B_k(x):=x^k(1-x)^k,
\qquad x\in\mathbb R.
\]
Expanding \((1-x)^k\) and differentiating \(k\) times gives Rodrigues' identity:
\[
\mathcal B_k^{(k)}(x)
=\sum_{m=0}^k(-1)^m\binom{k}{m}
  \frac{(k+m)!}{m!}x^m
=k!R_k(x).
\]
For every \(0\le i<k\), the factors \(x^k\) and \((1-x)^k\), respectively, ensure that
\[
\mathcal B_k^{(i)}(0)=\mathcal B_k^{(i)}(1)=0.
\]
Indeed, in each term of the Leibniz expansion, fewer than \(k\) derivatives fall on the factor vanishing at the relevant endpoint. Consequently, repeated integration by parts gives, for every polynomial \(R_j\),
\[
\int_0^1R_j(x)\mathcal B_k^{(k)}(x)\,dx
=(-1)^k\int_0^1R_j^{(k)}(x)\mathcal B_k(x)\,dx.
\]
For \(k=0\), this identity involves zero integrations by parts and has no boundary conditions.
% lean: Causalean.Mathlib.Analysis.SpecialFunctions.Jacobi.integral_mul_iterated_deriv_eq

First consider \(j=k\). The coefficient of \(x^k\) in \(R_k\) gives
\[
R_k^{(k)}(x)=k!(-1)^k\binom{2k}{k}.
\]
Combining this identity with Rodrigues' formula and integration by parts yields
\[
\begin{aligned}
k!\int_0^1R_k(x)^2\,dx
&=\int_0^1R_k(x)\mathcal B_k^{(k)}(x)\,dx\\
&=(-1)^k\int_0^1R_k^{(k)}(x)\mathcal B_k(x)\,dx\\
&=k!\binom{2k}{k}\int_0^1x^k(1-x)^k\,dx.
\end{aligned}
\]
The beta integral evaluates to
\[
\int_0^1x^k(1-x)^k\,dx
=\frac{k!}{(k+1)\prod_{i=0}^{k-1}(k+2+i)}
=\frac{(k!)^2}{(2k+1)!},
\]
where the product is \(1\) when \(k=0\). Since \(k!>0\), cancellation and the factorial formula for the binomial coefficient give
\[
\int_0^1R_k(x)^2\,dx
=\binom{2k}{k}\frac{(k!)^2}{(2k+1)!}
=\frac{1}{2k+1}.
\]
% lean: Causalean.Mathlib.Analysis.SpecialFunctions.Jacobi.shiftedLegendre_sq_integral_unit

If \(j<k\), the degree of \(R_j\) is \(j\), so \(R_j^{(k)}=0\). Thus
\[
k!\int_0^1R_j(x)R_k(x)\,dx
=(-1)^k\int_0^1R_j^{(k)}(x)\mathcal B_k(x)\,dx
=0,
\]
and the inner product vanishes. If \(k<j\), interchange the two factors and apply the same calculation. Hence
\[
\int_0^1R_j(x)R_k(x)\,dx
=
\begin{cases}
\dfrac{1}{2j+1},&j=k,\\
0,&j\ne k.
\end{cases}
\]
Multiplying the affine-substitution identity by \(2\) proves the asserted inner product on \([-1,1]\).
% lean: CausalSmith.Stat.RdTruesideNoiseFrontier.legendreP_orthogonal_exact

\item \emph{Reflection symmetry.}
The shifted Legendre reflection identity is
\[
R_j(x)=(-1)^jR_j(1-x),
\qquad x\in\mathbb R.
\]
Applying it at \(x=(1+t)/2\) gives
\[
P_j^{\mathrm{Leg}}(-t)
=R_j\!\left(\frac{1+t}{2}\right)
=(-1)^jR_j\!\left(\frac{1-t}{2}\right)
=(-1)^jP_j^{\mathrm{Leg}}(t),
\qquad t\in\mathbb R.
\]
% lean: CausalSmith.Stat.RdTruesideNoiseFrontier.legendreP_reflection_exact

\item \emph{The bound on the interval.}
The finite sum defining \(R_j\) is also the normalized shifted Jacobi polynomial with both shape parameters equal to zero: its coefficient identity is
\[
(-1)^m\binom jm
\frac{(j+m)!/j!}{m!}
=(-1)^m\binom jm\binom{j+m}{m},
\qquad 0\le m\le j.
\]
For \(\xi\in\mathbb R\) and \(r\in\mathbb N\), define the rising and falling factorial products by
\[
(\xi)_{\uparrow r}:=\prod_{i=0}^{r-1}(\xi+i),
\qquad
(\xi)_{\downarrow r}:=\prod_{i=0}^{r-1}(\xi-i),
\]
with both empty products equal to \(1\). Their definitions give
\[
(-\xi)_{\downarrow r}=(-1)^r(\xi)_{\uparrow r},
\qquad
(\xi)_{\uparrow m}(\xi+m)_{\uparrow(j-m)}
=(\xi)_{\uparrow j},
\quad 0\le m\le j.
\]
The falling-factorial convolution is
\[
(\xi+\eta)_{\downarrow j}
=\sum_{m=0}^j\binom jm
  (\xi)_{\downarrow m}(\eta)_{\downarrow(j-m)},
\qquad \xi,\eta\in\mathbb R.
\]
For nonnegative integer arguments, comparing coefficients of degree \(j\) in
\[
(1+\zeta)^{\xi+\eta}
=(1+\zeta)^\xi(1+\zeta)^\eta,
\qquad \xi,\eta\in\mathbb N,\quad \zeta\in\mathbb R,
\]
and multiplying by \(j!\) proves the convolution. Both sides are polynomials in the two arguments, so the identity extends to all real arguments.

Specialize the convolution to \(\xi=-(j+1)\) and \(\eta=j\). The product identities yield, for \(0\le m\le j\),
\[
\begin{aligned}
(-1)_{\downarrow j}&=(-1)^j j!,\\
(-(j+1))_{\downarrow m}
&=(-1)^m(j+1)_{\uparrow m}
=(-1)^m\frac{(j+m)!}{j!},\\
(j)_{\downarrow(j-m)}
&=(m+1)_{\uparrow(j-m)}
=\frac{j!}{m!}.
\end{aligned}
\]
Consequently the convolution becomes
\[
(-1)^j j!
=\sum_{m=0}^j\binom jm(-1)^m
  \frac{(j+m)!}{j!}\frac{j!}{m!}
=\sum_{m=0}^j(-1)^m\binom jm\frac{(j+m)!}{m!}.
\]
Dividing by \(j!>0\) supplies the normalization:
\[
R_j(1)
=\sum_{m=0}^j(-1)^m\binom jm\binom{j+m}{m}
=\frac{1}{j!}\sum_{m=0}^j(-1)^m\binom jm\frac{(j+m)!}{m!}
=(-1)^j,
\qquad |R_j(1)|=1.
\]
These identities include \(j=0\), when the products are empty and \(0!=1\).
At the other endpoint, direct evaluation of the defining sum gives
\[
R_j(0)=1,
\qquad |R_j(0)|=1.
\]
% lean: Causalean.Mathlib.Analysis.SpecialFunctions.Jacobi.jacobiShifted_one

When \(j=0\), the defining sum gives \(R_0(x)=1\) for every real \(x\), proving the required bound. For \(j\ge1\), define
\[
\lambda_j:=j(j+1)>0,
\qquad
\mathcal E_j(x):=
R_j(x)^2+\frac{x(1-x)R_j'(x)^2}{\lambda_j},
\qquad x\in\mathbb R.
\]
We derive the differential identity governing this energy. Define the coefficients over their full range by
\[
c_{j,m}:=
\begin{cases}
(-1)^m\binom jm\binom{j+m}{m},&0\le m\le j,\\
0,&m>j,
\end{cases}
\qquad m\in\mathbb N.
\]
The binomial coefficient formulas imply
\[
(m+1)^2c_{j,m+1}
=-(j-m)(j+m+1)c_{j,m},
\qquad 0\le m<j.
\]
The coefficient of \(x^m\) in the polynomial
\[
x(1-x)R_j''(x)+(1-2x)R_j'(x)+\lambda_jR_j(x)
\]
is
\[
(m+1)^2c_{j,m+1}
+\bigl(\lambda_j-m(m+1)\bigr)c_{j,m}.
\]
It vanishes for \(m<j\) by the recurrence, for \(m=j\) because \(c_{j,j+1}=0\) and \(\lambda_j=j(j+1)\), and for \(m>j\) because both coefficients are zero. Therefore
\[
x(1-x)R_j''(x)+(1-2x)R_j'(x)+\lambda_jR_j(x)=0,
\qquad x\in\mathbb R.
\]
% lean: Causalean.Mathlib.Analysis.SpecialFunctions.Jacobi.jacobiShifted_ode

Differentiation of the energy, followed by substitution of this differential equation, yields
\[
\begin{aligned}
\mathcal E_j'(x)
&=2R_j(x)R_j'(x)
+\frac{(1-2x)R_j'(x)^2
       +2x(1-x)R_j'(x)R_j''(x)}{\lambda_j}\\
&=(2x-1)\frac{R_j'(x)^2}{\lambda_j}.
\end{aligned}
\]
Thus the energy is nonincreasing on \([0,1/2]\) and nondecreasing on \([1/2,1]\). Its endpoint values are
\[
\mathcal E_j(0)=R_j(0)^2=1,
\qquad
\mathcal E_j(1)=R_j(1)^2=1.
\]
For \(x\le1/2\), compare with the left endpoint; for \(x>1/2\), compare with the right endpoint. These comparisons give
\[
\mathcal E_j(x)\le
\begin{cases}
\mathcal E_j(0),&0\le x\le1/2,\\
\mathcal E_j(1),&1/2<x\le1,
\end{cases}
\qquad
\mathcal E_j(x)\le1.
\]
Since the additional term in the energy is nonnegative on the unit interval,
\[
R_j(x)^2
\le R_j(x)^2+\frac{x(1-x)R_j'(x)^2}{\lambda_j}
=\mathcal E_j(x)\le1,
\qquad 0\le x\le1.
\]
Hence \(|R_j(x)|\le1\) throughout that interval. Finally,
\[
-1\le t\le1
\quad\Longrightarrow\quad
0\le\frac{1-t}{2}\le1
\quad\Longrightarrow\quad
|P_j^{\mathrm{Leg}}(t)|
=\left|R_j\!\left(\frac{1-t}{2}\right)\right|\le1.
\]
% lean: CausalSmith.Stat.RdTruesideNoiseFrontier.legendreP_abs_le_one_exact

\item \emph{The endpoint value.}
Evaluation at the positive endpoint gives
\[
P_j^{\mathrm{Leg}}(1)=R_j(0)=1.
\]
Applying reflection symmetry at \(t=1\) therefore yields
\[
P_j^{\mathrm{Leg}}(-1)
=(-1)^jP_j^{\mathrm{Leg}}(1)
=(-1)^j.
\]
% lean: CausalSmith.Stat.RdTruesideNoiseFrontier.legendreP_neg_one_exact
\end{enumerate}
\end{proof}
